
% ============================================================
% File: VS.B01_ONTOLOGY_OF_VIABILITY/1_FRONTMATTER/09_structure_and_referencing_guide.tex
% Book: Ontology of Viability
% Purpose: Structure and Referencing Guide
% ============================================================

\cleardoublepage

\chapter*{Structure and Referencing Guide}
\addcontentsline{toc}{chapter}{Structure and Referencing Guide}

% ------------------------------------------------------------
% 1. PURPOSE
% ------------------------------------------------------------

\section*{1. Purpose of This Guide}

This guide explains the structural organization,
documentary principles, and referencing conventions
of the present volume, \textit{Ontology of Viability}.

The book is not organized merely as a sequence of
independent chapters. Its structure is intended
to preserve conceptual dependencies, argumentative
continuity, and the traceability of the philosophical
investigation.

The organization of the volume serves several
related purposes:

\begin{itemize}
    \item To distinguish the major philosophical domains
    investigated in the book.
    \item To preserve relationships between parts,
    chapters, sections, and subordinate textual units.
    \item To make the development of foundational concepts
    traceable across the volume.
    \item To distinguish foundational commitments
    from provisional constructions, hypotheses,
    and unresolved questions.
    \item To maintain consistency between the main text,
    supporting registers, bibliography, glossary,
    and other reference materials.
    \item To facilitate justified revision without
    obscuring the development of the philosophical
    argument.
\end{itemize}

The architecture of the volume is therefore both
editorial and methodological. It organizes the
investigation while supporting the systematic
development and examination of its philosophical
content.

The documentary organization must not, however,
be confused with the ontological structure of
reality itself. The former provides a means of
presenting and tracing the investigation; the
latter must be examined and justified through
the philosophical arguments developed in the book.

% ------------------------------------------------------------
% 2. OVERALL ARCHITECTURE
% ------------------------------------------------------------

\section*{2. The Overall Architecture of the Volume}

The volume is organized into three principal
documentary layers: front matter, main matter,
and back matter.

These layers have different functions and should
not be treated as interchangeable.

\subsection*{2.1. Front Matter}

The front matter introduces the book, establishes
its initial orientation, and provides the reader
with the information required to understand
its purpose, scope, and methodological approach.

Its principal functions include:

\begin{itemize}
    \item Bibliographic and editorial identification.
    \item Presentation of the intellectual purpose
    of the volume.
    \item Introduction of the central ontological
    problems.
    \item Explanation of the methodological orientation.
    \item Guidance for reading and interpreting the volume.
    \item Explanation of the documentary and referencing
    conventions.
\end{itemize}

The front matter does not replace the philosophical
demonstrations developed in the main matter.
Its function is to establish the orientation
within which those demonstrations can be examined.

\subsection*{2.2. Main Matter}

The main matter contains the principal philosophical
investigation.

It is divided into four parts:

\begin{description}

    \item[Part I: Identity and Manifestation]
    This part establishes the initial ontological
    problems concerning Being, identity, determination,
    unity, distinction, individuation, numerical identity,
    continuity, and manifestation.

    Its central commitment is the ontological
    inseparability of Being and identity, without
    assigning priority to either.

    \item[Part II: The Realms of Being]
    This part investigates the differentiation
    of Being through a framework comprising four
    fundamental realms and three intermediate domains.

    The four fundamental realms are First Determination,
    Structural Determination, Relational Manifestation,
    and Material Manifestation.

    The three intermediate domains are Hub of Decree,
    Hub of Signification, and Hub of Realization.

    These domains must not be interpreted as seven
    equivalent levels arranged in a simple linear
    hierarchy. Their precise functions and relations
    must be examined through the arguments developed
    in the relevant chapters.

    The possible status and location of a Transcendent
    Ground or Throne-Domain remain open unless
    established through sufficient ontological argument.

    \item[Part III: Light]
    This part investigates Light as an ontological
    problem.

    Its meaning and function must be established
    within the ontological architecture of the volume
    rather than reduced to a metaphor for human
    understanding or the acquisition of knowledge.

    \item[Part IV: Dynamics of Being]
    This part investigates the dynamic dimension
    of Being, including becoming, transformation,
    interaction, continuity, discontinuity,
    emergence, temporality, generation, dissolution,
    reversibility, irreversibility, cost, residue,
    and consequence.

    Its scope extends beyond observable movement
    or empirical change to the ontological conditions
    and implications of transformation.

\end{description}

The four parts are distinct but interconnected.
Their order represents the architecture of the
philosophical exposition and should not automatically
be interpreted as a temporal sequence or a complete
hierarchy of ontological dependence.

The inclusion of a topic within a particular part
does not by itself establish that every proposition
associated with that topic has already been fully
defined or demonstrated.

\subsection*{2.3. Back Matter}

The back matter provides supplementary materials
that support the interpretation, examination,
and continued development of the volume.

Its principal components include:

\begin{description}

    \item[Conclusion:]
    Provides a synthesis of the ontological
    investigation and identifies its principal
    results, limitations, and unresolved questions.

    \item[Appendices:]
    Contain supplementary materials that are necessary
    or useful for the investigation but do not belong
    in the continuous development of the main chapters.

    \item[Logical Entity Reference Map:]
    Provides a structured representation of
    relationships among identified philosophical
    entities and their documentary locations.

    \item[Bibliography:]
    Organizes the sources cited or examined
    in the volume.

    \item[Glossary:]
    Provides controlled explanations of the principal
    terms used throughout the text.

    \item[Indexes:]
    Facilitate the retrieval of concepts, names,
    subjects, and other relevant elements.

    \item[Registers:]
    Provide structured records of philosophical
    entities and argumentative elements identified
    in the investigation.

\end{description}

The back matter supplements the main argument.
It must not be used to introduce unsupported
conclusions that should instead be established
within the relevant philosophical chapters.

% ------------------------------------------------------------
% 3. STRUCTURAL AND CONCEPTUAL TRACEABILITY
% ------------------------------------------------------------

\section*{3. Structural and Conceptual Traceability}

An important requirement of the volume is that
its philosophical claims remain traceable across
their different textual and documentary locations.

Traceability means that a reader or editor should
be able to identify:

\begin{itemize}
    \item Where a concept is introduced.
    \item Where it is formally defined.
    \item Where it is examined, qualified, or revised.
    \item Which arguments support or challenge it.
    \item Whether it is treated as a foundational
    commitment, provisional construction, hypothesis,
    or unresolved question.
    \item Where it is recorded in the relevant register.
    \item Where it is connected to other concepts
    and philosophical claims.
\end{itemize}

This requirement is particularly important because
the volume develops an interconnected ontological
framework rather than a collection of independent
definitions.

\subsection*{3.1. Conceptual Identity and Documentary Identity}

A philosophical concept and its documentary record
are not the same kind of entity.

The concept belongs to the philosophical content
of the investigation. Its documentary record is
a structured representation that allows the concept
to be identified, retrieved, and examined.

A stable identifier can therefore be assigned
to a conceptual entity without implying that
the identifier constitutes or determines its
ontological identity.

The identifier serves an editorial and documentary
function. It does not establish the existence,
nature, or ontological independence of the entity
to which it refers.

\subsection*{3.2. Stable Identifiers}

Where stable identifiers are adopted, they should
satisfy the following requirements:

\begin{itemize}
    \item Uniqueness within the relevant category
    or register.
    \item Persistence across revisions.
    \item Consistency between the main text and
    documentary records.
    \item Independence from temporary page numbers.
    \item Clarity concerning the category of the
    recorded entity.
\end{itemize}

Identifiers should not be reassigned merely because
a paragraph, section, or chapter has moved during
revision.

If the conceptual content of an entity changes
substantially, the change should be documented
rather than concealed through an unexplained
identifier replacement.

The documentary continuity of an identifier must
not be confused with the philosophical continuity
or numerical identity of the concept it represents.

% ------------------------------------------------------------
% 4. CONCEPTUAL CATEGORIES AND THEIR DOCUMENTARY
%    FUNCTIONS
% ------------------------------------------------------------

\section*{4. Conceptual Categories and Their Documentary Functions}

The referencing system must preserve distinctions
between different types of philosophical content.

These categories organize the investigation.
They do not independently determine the truth
or ontological status of the claims they contain.

\subsection*{4.1. Axioms}

An axiom records a foundational commitment adopted
within the philosophical framework.

Its status must be explicit. It should not be
confused with an ordinary observation, a provisional
interpretation, or an unexamined assumption.

\subsection*{4.2. Definitions}

A definition establishes how a term or concept
is used within the system.

Definitions should be sufficiently precise to
support consistent reasoning and distinguish
the intended concept from related or potentially
ambiguous concepts.

A definition establishes terminological usage;
it does not automatically establish the independent
ontological validity of what it defines.

\subsection*{4.3. Propositions}

A proposition expresses a claim that can be
examined through logical, conceptual, or
argumentative analysis.

Its relationship to premises, definitions,
and other propositions should be made explicit
whenever relevant.

\subsection*{4.4. Theses}

A thesis expresses a principal philosophical
position advanced by the investigation.

A thesis may be supported by several arguments
and may depend on multiple propositions.

Its status should reflect the strength and
completeness of the justification developed
in the text.

\subsection*{4.5. Hypotheses}

A hypothesis is a provisional explanatory
or ontological proposal that requires further
examination.

Its provisional status must remain visible
in both the main text and the relevant register.

A hypothesis should not be silently transformed
into an established commitment.

\subsection*{4.6. Arguments}

An argument records a structured relation
between premises, inferential steps, and
a conclusion.

Where possible, its documentary record should
identify the principal premises, the conclusion,
the relevant chapter, and important dependencies.

An argument should be evaluated according to
its actual structure and inferential adequacy,
not merely according to the philosophical
importance of its conclusion.

\subsection*{4.7. Philosophical Foundations}

Philosophical foundations record deeper
commitments, conceptual orientations, and
foundational considerations that support
the development of the system.

They may overlap in subject matter with axioms,
theses, or definitions, but their documentary
function is distinct.

The relevant records should preserve these
distinctions rather than forcing every foundational
statement into a single category.

% ------------------------------------------------------------
% 5. CROSS-REFERENCING
% ------------------------------------------------------------

\section*{5. Cross-Referencing Within the Main Text}

Cross-references allow the reader to move between
related parts of the investigation without treating
each chapter as conceptually isolated.

The following principles should guide their use:

\begin{itemize}
    \item References to chapters should identify
    the relevant chapter rather than relying
    exclusively on page numbers.
    \item References to sections should use
    the document's actual section labels.
    \item References to propositions, definitions,
    or arguments should use stable identifiers
    when such identifiers have been formally adopted.
    \item References to earlier discussions should
    clarify whether the earlier material is being
    recalled, extended, qualified, or challenged.
    \item References to later discussions should
    not be used as substitutes for arguments
    necessary in the current context.
\end{itemize}

LaTeX labels and references should be used wherever
appropriate to maintain consistency after document
revisions.

For example, a chapter or section may be referenced
through the standard LaTeX mechanisms:

\begin{verbatim}
\label{...}
\ref{...}
\pageref{...}
\end{verbatim}

The actual label names should follow a consistent
naming convention.

These examples describe the referencing mechanism
and do not establish specific label names for
chapters or sections that have not yet been defined.

% ------------------------------------------------------------
% 6. REFERENCES TO REGISTERS
% ------------------------------------------------------------

\section*{6. References to Registers}

The registers provide a documentary layer that
complements the continuous philosophical exposition.

A reference to a register entry should be used
when the reader needs to inspect the structured
record of a philosophical entity or claim.

Such a reference is particularly useful when:

\begin{itemize}
    \item A concept is used repeatedly across
    different chapters.
    \item A proposition depends on several
    other propositions.
    \item A thesis has a complex argumentative history.
    \item A hypothesis is revised or replaced.
    \item A foundational commitment is connected
    to multiple philosophical consequences.
    \item The status of a claim changes during
    the development of the volume.
\end{itemize}

The main text should remain intelligible without
requiring the reader to consult a register after
every paragraph.

Registers are intended to improve transparency
and traceability, not to fragment the philosophical
argument.

\subsection*{6.1. Registry Consistency}

The following consistency requirements should
be maintained:

\begin{itemize}
    \item An identifier must refer to one identifiable
    record within its designated category.
    \item A record must not be assigned conflicting
    statuses in different documentary locations.
    \item A revised claim should be reflected
    in its relevant register entry.
    \item References to deleted or superseded
    records should be resolved or explicitly
    documented.
    \item The main text and its supporting records
    should not silently contradict one another.
\end{itemize}

Where a claim is revised, its documentary history
should be preserved whenever that history is
philosophically or methodologically significant.

% ------------------------------------------------------------
% 7. LOGICAL ENTITY REFERENCE MAP
% ------------------------------------------------------------

\section*{7. Logical Entity Reference Map}

The Logical Entity Reference Map supports
the identification of relationships among
philosophical entities.

Its purpose is broader than listing terms
or pointing to individual pages.

It should help establish documentary connections
between:

\begin{itemize}
    \item Concepts and their definitions.
    \item Axioms and dependent propositions.
    \item Premises and argumentative conclusions.
    \item Theses and their supporting arguments.
    \item Hypotheses and the questions they address.
    \item Foundational commitments and the chapters
    in which they are developed.
    \item Concepts that are related, distinguished,
    or placed in tension with one another.
\end{itemize}

The reference map should distinguish different
types of relationships.

Conceptual dependence, argumentative support,
terminological equivalence, contrast, and
unresolved tension should not be represented
as if they were the same relation.

The map is a documentary and analytical instrument.
It does not establish the ontological reality
of a relationship merely by representing it.

A represented connection must remain distinguishable
from a philosophically demonstrated connection.

% ------------------------------------------------------------
% 8. BIBLIOGRAPHIC REFERENCING
% ------------------------------------------------------------

\section*{8. Bibliographic Referencing}

All external sources used in the investigation
should be referenced in a consistent and
verifiable manner.

Bibliographic references serve several purposes:

\begin{itemize}
    \item Identifying the intellectual origin
    of cited material.
    \item Distinguishing direct quotation from
    paraphrase.
    \item Enabling readers to locate the original source.
    \item Clarifying the relationship between
    the present argument and previous philosophical
    discussions.
    \item Supporting comparative analysis without
    obscuring the independent structure of the
    present framework.
\end{itemize}

A reference to an external source must not be
treated as a substitute for philosophical
justification.

The fact that a position has been advanced
by an influential thinker does not, by itself,
establish its validity.

Similarly, disagreement with an existing
philosophical position should be directed
toward its actual conceptual or argumentative
content.

\subsection*{8.1. Primary Sources}

Primary sources should be cited when the
investigation engages directly with an original
philosophical text or another original source.

Where interpretation depends on a particular
passage, the relevant edition and location
should be identified as precisely as possible.

\subsection*{8.2. Secondary Sources}

Secondary sources should be used to examine
interpretations, scholarly discussions,
historical developments, and critical analyses.

They should not be presented as if they were
the original statements of the thinkers or
authors being discussed.

\subsection*{8.3. Comparative Sources}

Comparative sources support the examination
of similarities, differences, conceptual
correspondences, and disagreements between
the present framework and other philosophical
or intellectual traditions.

Comparison should not automatically imply
equivalence.

A shared term may have different conceptual
functions in different systems. A meaningful
comparison must therefore examine both
terminology and argumentative context.

\subsection*{8.4. Citation Integrity}

Every quotation, paraphrase, and source-dependent
claim should be represented accurately.

The following principles apply:

\begin{itemize}
    \item Quotations must preserve the wording
    and context of the source.
    \item Paraphrases must not distort the
    original position.
    \item Interpretations should be distinguished
    from direct statements made by the source.
    \item Bibliographic records should contain
    sufficient information for source identification.
    \item Uncertain publication details must
    not be invented.
    \item Sources that have not been consulted
    directly should not be represented as directly
    examined.
\end{itemize}

% ------------------------------------------------------------
% 9. GLOSSARY AND TERMINOLOGICAL CONSISTENCY
% ------------------------------------------------------------

\section*{9. Glossary and Terminological Consistency}

The glossary provides a controlled reference
for the principal terms used in the volume.

Its purpose is not merely to translate or
summarize terminology. It should preserve
distinctions that are important to the
philosophical structure.

A term may have a familiar meaning in ordinary
language or in another philosophical tradition
while acquiring a more specific function within
the present framework.

Where such differences occur, the glossary
should clarify the intended usage.

\subsection*{9.1. Definitions and Contextual Usage}

A glossary entry should be consistent with
the formal definition of a concept when such
a definition exists.

However, the glossary should not be treated
as the sole location in which the meaning
of a central concept is established.

The main text remains responsible for developing
and justifying the conceptual role of the term.

\subsection*{9.2. Terminological Revision}

If the meaning of a term changes during
the development of the system, the revision
should be documented.

A term should not retain an unchanged glossary
definition while being used in a substantially
different sense in the main text.

Where a distinction between earlier and later
usage is necessary, the relevant conceptual
development should be made explicit.

% ------------------------------------------------------------
% 10. INDEXES AND RETRIEVAL
% ------------------------------------------------------------

\section*{10. Indexes and Retrieval}

Indexes support the retrieval of important
subjects and conceptual locations.

They should help readers locate:

\begin{itemize}
    \item Central ontological concepts.
    \item Principal conceptual distinctions.
    \item Philosophical positions.
    \item Significant arguments.
    \item Relevant thinkers and sources.
    \item Important relations among concepts.
\end{itemize}

An index is a retrieval instrument rather
than a substitute for the conceptual structure
of the book.

The occurrence of a term in several locations
does not necessarily mean that it performs
the same philosophical function in every
occurrence.

Readers should therefore interpret indexed
locations within their respective argumentative
contexts.

% ------------------------------------------------------------
% 11. FIGURES AND FORMAL REPRESENTATIONS
% ------------------------------------------------------------

\section*{11. Figures, Diagrams, and Formal Representations}

Visual and formal materials may be used
to clarify the relationships developed
in the main text.

Their function is supplementary and explanatory.

A figure or diagram should be accompanied
by sufficient textual explanation to establish:

\begin{itemize}
    \item What entities or concepts it represents.
    \item What relationships it intends to show.
    \item What its boundaries and limitations are.
    \item Whether it represents an established
    structure, a provisional model, or a
    hypothetical arrangement.
\end{itemize}

A visual representation must not imply
a degree of certainty or completeness that
the corresponding argument has not established.

In particular, the representation of the
seven domains must preserve the distinction
between the four fundamental realms and
the three intermediate domains.

It must not inadvertently reduce them to
seven equivalent linear levels.

Formal representations should be interpreted
as expressions of the conceptual framework
under examination. They do not independently
establish the ontological validity of that
framework.

% ------------------------------------------------------------
% 12. REVISION AND VERSION CONSISTENCY
% ------------------------------------------------------------

\section*{12. Revision and Version Consistency}

The volume is a developing philosophical
construction. Its documentary organization
must therefore support revision without
sacrificing conceptual continuity.

A revision may affect:

\begin{itemize}
    \item The wording or scope of a definition.
    \item The formulation of a proposition.
    \item The argumentative structure of a thesis.
    \item The status of a hypothesis.
    \item The relation between two concepts.
    \item The organization of a chapter.
    \item The interpretation of a foundational
    commitment.
\end{itemize}

Such changes should be reflected in the
relevant documentary locations.

The following principles should guide revision:

\begin{itemize}
    \item Preserve stable identifiers where
    the underlying conceptual entity remains
    identifiable.
    \item Document substantial changes in
    conceptual meaning.
    \item Update cross-references affected
    by structural changes.
    \item Maintain consistency between the
    main text and supporting registers.
    \item Distinguish corrections of expression
    from changes in philosophical position.
    \item Avoid presenting provisional revisions
    as established conclusions without adequate
    justification.
\end{itemize}

The purpose of version consistency is not
to prevent philosophical development.
It is to make that development intelligible
and traceable.

% ------------------------------------------------------------
% 13. STRUCTURE AND PHILOSOPHICAL CONTENT
% ------------------------------------------------------------

\section*{13. The Relationship Between Structure and Philosophical Content}

The architecture of this volume is designed
to support the philosophical investigation,
not to predetermine its conclusions.

The existence of a chapter, register, diagram,
or category does not by itself establish
the ontological validity of the subject
it contains.

Similarly, the documentary distinction
between two entities does not necessarily
imply that they are ontologically independent.

This distinction is particularly important
in the present investigation, where conceptual
differentiation must remain compatible with
ontological inseparability wherever the
arguments require it.

The distinction between Being and identity,
for example, is analytically necessary but
must not be interpreted as establishing
an ontological priority of either term.

Likewise, the fact that the volume discusses
the relation between ontological reality
and its manifestation must not be interpreted
as making Being dependent on its documentary
representation.

The structure of the book should therefore
be read as an organizational representation
of the investigation rather than as a complete
ontological map established independently
of its arguments.

% ------------------------------------------------------------
% 14. STATUS OF PHILOSOPHICAL CLAIMS
% ------------------------------------------------------------

\section*{14. The Status of Philosophical Claims}

The referencing and documentary architecture
must preserve the philosophical status of
each recorded claim.

The following distinctions are especially
important:

\begin{description}

    \item[Foundational commitments:]
    Positions adopted as constitutive starting
    points of the present ontological framework.

    \item[Provisional constructions:]
    Conceptual or structural arrangements
    that are being developed and whose final
    form may remain subject to revision.

    \item[Hypotheses:]
    Explanatory or ontological proposals
    that require further examination.

    \item[Open questions:]
    Problems for which the investigation
    has not yet established a sufficient answer.

    \item[Rejected or superseded positions:]
    Earlier formulations that have been
    explicitly revised, qualified, or abandoned.

\end{description}

These categories should not be collapsed
into a single undifferentiated record of
philosophical claims.

The unresolved status of a question is not
necessarily a defect in the investigation.

An explicit record of what remains unresolved
can be more methodologically valuable than
an unsupported attempt to produce premature
closure.

% ------------------------------------------------------------
% 15. PRACTICAL PRINCIPLES
% ------------------------------------------------------------

\section*{15. Practical Principles for Readers and Editors}

For practical use, the following principles
summarize the intended operation of the
documentary system:

\begin{enumerate}
    \item Read the main text as a connected
    philosophical investigation rather than
    as a collection of independent essays.

    \item Use the part and chapter structure
    to understand the progression of the
    ontological inquiry.

    \item Consult the glossary when a term
    has a specific technical meaning within
    the framework.

    \item Consult the relevant register when
    the status, identifier, or documentary
    history of a claim requires clarification.

    \item Use cross-references to follow
    conceptual and argumentative dependencies.

    \item Consult the Logical Entity Reference
    Map when relationships among multiple
    philosophical entities need to be examined.

    \item Use the bibliography to distinguish
    primary engagement, secondary interpretation,
    and comparative analysis.

    \item Interpret figures and diagrams
    in conjunction with their textual
    explanations.

    \item Treat hypotheses and open questions
    according to their explicitly stated status.

    \item Evaluate philosophical claims through
    their arguments rather than through their
    documentary classification alone.
\end{enumerate}

% ------------------------------------------------------------
% 16. CONCLUDING STATEMENT
% ------------------------------------------------------------

\section*{16. Concluding Statement}

The documentary architecture of
\textit{Ontology of Viability} is intended
to make the development of its philosophical
content systematic, transparent, and traceable.

Its organizational structure distinguishes
the major domains of investigation while
preserving their conceptual interdependence.

Its registers distinguish different categories
of philosophical claims without confusing
documentation with justification.

Its referencing system connects concepts,
definitions, arguments, propositions, theses,
hypotheses, and sources across the volume.

Its revision principles allow the framework
to develop without concealing the history
or status of its conceptual constructions.

Ultimately, the architecture of the book
serves a methodological purpose: to support
a rigorous investigation of Being while
preventing the organization of the philosophical
account from being mistaken for the totality
of reality itself.

The documentary system must remain answerable
to the philosophical investigation. It should
clarify the structure of the inquiry, preserve
its distinctions, and make its development
intelligible, while leaving the ontological
validity of substantive claims to be established
through the appropriate philosophical arguments.

\cleardoublepage

% ============================================================
% End of Structure and Referencing Guide
% ============================================================
